\section{Discussion}
\label{sec:discussion}

\paragraph{Two readable signals and one causal gate.}
Our results separate three things that the phrase ``an internal state transition'' tends to conflate. \emph{Sufficiency} is a property of the input: whether the context now contains every fact the question needs. The model represents it cleanly, but a text classifier that sees the same context can compute it too, and its causal counterpart in the model is an abstention gate. \emph{Uptake} is a property of the model: whether the same decisive paragraph will actually be used in the answer. It is predictable from activations but not from text, at $0.78$ AUROC across three model families and on every external source. Neither readout direction is a lever that controls the behaviour it predicts. What turns an insufficient state into an answer is high-dimensional, orthogonal to both readout directions, and on some models concentrated on the answer's unembedding direction. The picture is therefore a low-dimensional gate that decides whether to answer, plus high-dimensional content that determines the answer, rather than a single direction along which the model takes up new evidence.

\paragraph{Why probes do not find the lever, and what to use instead.}
A logistic probe looks for the direction along which the two classes are best separated relative to their spread. It therefore down-weights directions along which activations vary a lot within each class. The difference-of-means direction, in contrast, follows the raw shift between the class averages, and much of that shift lies along a high-variance axis related to answering. The projection-out and subspace experiments show that the causal content lies exactly there. A probe direction can thus score a high AUROC and transfer well across datasets while being causally irrelevant. Any proposed steering direction should therefore be validated with controls of the kind used here: projecting the direction out, measuring how many dimensions the effect needs, and comparing with random directions of matched norm. For retrieval systems, the most useful signal is uptake, as a detector of evidence the model has ignored. Sufficiency is tracked about as well by the model's own abstentions, and the abstention gate trades abstention against compliance uniformly rather than making the model more sensitive to the evidence.
